\documentclass[10pt,a4paper]{amsart}
\usepackage[margin=19mm]{geometry}
\usepackage{amsmath,amssymb,mathtools}
\usepackage[scr=rsfs,cal=euler]{mathalfa}
\usepackage{xfrac}
\usepackage[unicode,hidelinks,bookmarks=false]{hyperref}
\newcommand{\ie}{i.e.\ }
\newcommand{\cf}{cf.\ }
\newcommand{\resp}{respectively\ }
\newcommand{\Subsection}{\subsection}
\providecommand{\nobreakdash}{\nobreak\hbox{-}}
\setlength{\emergencystretch}{2em}
\multlinegap0pt
\usepackage{amssymb}
\usepackage{mathtools}
\usepackage{tikz}
\newtheorem{lemma}{Lemma}
\title{On prime endomorphisms of the free group of rank two}
\author{Andreas Thom}
\date{Source: arXiv:2606.02225v1 (CC BY 4.0)}
\begin{document}
\maketitle

\noindent\emph{Source and licence.} On prime endomorphisms of the free group of rank two. Version arXiv:2606.02225v1 (CC BY 4.0), distributed by arXiv under the Creative Commons Attribution 4.0 International licence, http://creativecommons.org/licenses/by/4.0/, as recorded in the arXiv licence metadata for this exact version and shown on the abstract page https://arxiv.org/abs/2606.02225v1. Full licence notice: \texttt{licenses/arxiv-2606.02225-CC-BY-4.0.txt}.\par\smallskip

\noindent\textit{Editorial scope.} Counted unit: the article's lemma L.return (source0.tex lines 906--914). The article's complete original proof of it is delivered with it (source0.tex lines 916--1005, one \texttt{proof} environment of 90 lines). No mathematics is written, repaired or re-derived: the statement and the proof are the article's own lines, in the article's own notation, and no retained line is substituted or re-flowed.  Retained uncounted as the source-local context the proof needs: lines 836--848.  Nothing was omitted from any retained span, so the package carries no omission record.  No retained line contains a citation, so the wrapper carries no bibliography and no bibliography entry.  No retained line contains a figure environment, an image-inclusion command or drawing code, so no image is delivered and the package declares no asset dependency.  Editorial formatting only, in addition: an AMS \texttt{amsart} preamble that restates the article's own declarations and macros used by the retained lines, namely the environments \texttt{lemma} on separate counters, together with the standard packages those lines need.  The retained environments print in this document as Lemma 1 in order of appearance, whereas the article prints the same items with the numbering of its own full document.  The complete native source of the article as distributed in the arXiv e-print for this exact version is bundled unchanged as \texttt{sources/201886/source0.tex}.
\begin{lemma}\label{L.sharp-cogrowth}
Let
\[
        \rho_0=
        \frac{1+\sqrt[3]{46+6\sqrt{57}}+\sqrt[3]{46-6\sqrt{57}}}{3}.
\]
For every proper rank-two subgroup $K<F_2$, one has
\[
        \rho(\Gamma(K))\leq \rho_0.
\]
Moreover, equality can occur only when the topological core of $K$ is a
figure-eight with loop lengths $1$ and $2$.
\end{lemma}

\begin{lemma}\label{L.return}
There is a constant $C_0>0$ such that, for every proper rank-two subgroup
$K<F_2$, if $h(K)$ denotes the stem length of $\widehat\Gamma(K)$, then for
every $n\geq 1$,
\[
        \frac{|K\cap S_n|}{|S_n|}
        \leq C_0\left(\frac34\right)^n\left(\frac94\right)^{-2h(K)}.
\]
\end{lemma}

\begin{proof}
Let $X=\Gamma(K)$, let $h=h(K)$, and let $x\in X$ be the attachment point of
the stem of $\widehat\Gamma(K)$.  Let
\[
        F_{X,x}(z)=\sum_{n\geq 1} c_n(X,x)z^n
\]
be the generating series for non-backtracking closed paths in $X$ based at $x$.
By the preceding discussion, the generating series
\[
        F_K(z)=\sum_{n\geq 1}|K\cap S_n|z^n
\]
for reduced words in $K$ satisfies
\[
        F_K(z)=z^{2h}F_{X,x}(z).
\]
We therefore estimate $F_{X,x}(z)$ uniformly.

The point $x$ is either a branch vertex
of $\bar X$ or lies in the interior of a unique topological edge.  In
either case there are only boundedly many possible initial partial segments
leaving $x$ and terminal partial segments returning to $x$.  Consequently
there are vectors $u_{X,x}(z),v_{X,x}(z)\in \mathbb C^{\Omega_X}$ whose entries are
finite sums of at most uniformly many monomials in $z$, and a function $g_X(z)$
such that
\[
        F_{X,x}(z)=g_X(z)+u_{X,x}(z)^\ast\bigl(I-M_X(z)\bigr)^{-1}v_{X,x}(z).
\]
Here $g_X(z)$ is either a polynomial with coefficients bounded by an absolute
constant, or else $x$ lies in the interior of a topological loop of length
$\ell$ and
\[
        g_X(z)=2\sum_{n\geq 1} z^{n\ell}=\frac{2z^\ell}{1-z^\ell}.
\]
Indeed, if $x$ does not lie in the interior of a topological loop, then a
based non-backtracking closed path that never reaches a branch vertex has only
finitely many possibilities, and these are recorded by a polynomial.  If $x$
does lie in the interior of a topological loop of length $\ell$, then the closed
non-backtracking paths that stay entirely in that loop are exactly the nonzero
windings around the loop in one of the two orientations, giving the displayed
geometric series.  All remaining based closed paths are obtained by choosing an
initial partial segment from $x$ to a branch vertex, following a
non-backtracking itinerary of oriented topological edges, and then taking a
terminal partial segment back to $x$.  This is encoded by the displayed
formula.

Set
$
\rho_*=\frac94,
$ and
$       R=\rho_*^{-1}=\frac49.$
By Lemma~\ref{L.sharp-cogrowth}, every proper rank-two core has spectral radius
at most $\rho_0<\rho_*$.  Therefore, at $R=4/9$, the weighted transfer matrices
are uniformly inside their disk of convergence.  More explicitly, for
$|z|\leq R$, the entrywise absolute value of $M_X(z)$ is dominated by one of the
finitely many maximal matrices corresponding to the three topological types: the
figure-eight of lengths $(1,2)$, the theta graph of lengths $(1,1,1)$, or the
barbell graph of lengths $(1,1,1)$.  In the theta and barbell cases the maximal
matrix has row sums at most $2R=8/9<1$.  In the figure-eight case the maximal
matrix is
\[
        \begin{pmatrix}
                R & 2R^2 \\
                2R & R^2
        \end{pmatrix},
\]
whose spectral radius is strictly smaller than $1$ because its critical value is
$\rho_0^{-1}>R$.  Hence the Neumann series
\[
        \bigl(I-M_X(z)\bigr)^{-1}=\sum_{j\geq 0} M_X(z)^j
\]
is uniformly bounded for all proper rank-two $K$ and all $|z|\leq R$.
In the polynomial case $g_X(z)$ is uniformly bounded there, and in the loop case
one has
\[
        |g_X(z)|\leq \frac{2R}{1-R}=\frac{8}{5}.
\]
Since each entry of $u_{X,x}(z)$ and $v_{X,x}(z)$ is a sum of at most uniformly many
monomials $z^d$ and $|z|\leq R<1$, these entries are also uniformly bounded
there,
the generating series $F_{X,x}(z)$ is bounded by an absolute constant on that disk,
independently of $K$.

Therefore $F_K(z)=z^{2h}F_{X,x}(z)$ is bounded by $C_1R^{2h}$ on $|z|\leq R$.
By Cauchy's estimate, the coefficient of $z^n$ in $F_K(z)$ satisfies
\[
        |K\cap S_n|\leq C_1 R^{2h-n}=C_1 \rho_*^{\,n-2h}.
\]
Since $|S_n|=4\cdot 3^{n-1}$, we obtain
$\frac{|K\cap S_n|}{|S_n|}\leq C_0\left(\frac34\right)^n\left(\frac94\right)^{-2h}$.
\end{proof}

\end{document}
